\documentclass[12pt, letterpaper]{article}

% ── Packages ────────────────────────────────────────────────
\usepackage[margin=1in]{geometry}
\usepackage{setspace}
\usepackage{parskip}
\usepackage{titlesec}
\usepackage{graphicx}
\usepackage{amsmath, amssymb}
\usepackage{booktabs}
\usepackage{caption}
\usepackage{subcaption}
\usepackage{hyperref}
\usepackage{fancyhdr}
\usepackage{times}  % Times New Roman font
\usepackage{microtype}
\usepackage{enumitem}
\usepackage{float}  % [H] figure placement
\usepackage{multirow}

% ── Hyperlink Styling ────────────────────────────────────────
\hypersetup{
    colorlinks = true,
    linkcolor  = blue,
    citecolor  = blue,
    urlcolor   = blue,
}

% ── Section / Subsection Formatting ─────────────────────────
\titleformat{\section}{\large\bfseries}{}{0em}{}[\titlerule]
\titleformat{\subsection}{\normalsize\bfseries}{}{0em}{}
\titleformat{\subsubsection}{\normalsize\itshape}{}{0em}{}

% ── Header / Footer ─────────────────────────────────────────
\pagestyle{fancy}
\fancyhf{}
\rhead{\small CS 679 Pattern Recognition}
\lhead{\small Assignment 3}
\cfoot{\thepage}
\renewcommand{\headrulewidth}{0.4pt}

% ── Spacing ─────────────────────────────────────────────────
\onehalfspacing

% ============================================================
%  DOCUMENT
% ============================================================
\begin{document}

% ── COVER PAGE ──────────────────────────────────────────────
\begin{titlepage}
	\centering
	\vspace*{2cm}

	{\LARGE \textbf{CS 679 Pattern Recognition}}\\[0.5cm]

	{\Large \textbf{Assignment 3 - EigenFaces}}\\[0.5cm]

	\vfill

	{\large \textbf{Richard White}}\\[0.3cm]
	{\normalsize May 4, 2026}\\[1cm]

	\vfill

	% ── Academic Integrity Statement ────────────────────────
	\begin{minipage}{0.85\textwidth}
		\small
		\textit{``I declare that all material in this assignment is my own work except where there is clear
			acknowledgment or reference to the work of others. I understand that both my report and code may be
			subjected to plagiarism detection software, and fully accept all consequences if found responsible for
			plagiarism, as explained in the syllabus, and described in UNR’s Academic Standards Policy: UAM
			6,502.''}

		\vspace{1cm}
		\noindent\rule{6cm}{0.4pt}\\
		Richard White
	\end{minipage}

	\vspace{2cm}
\end{titlepage}

% ── TABLE OF CONTENTS ───
\tableofcontents
\newpage

% ============================================================
%  SECTION 1 -- THEORY
% ============================================================
\section{Theory}

\subsection{Face Space Construction}

% This section covers the foundational steps for building the face space model during the training phase.

% \paragraph{The Average Face}

Face space construction utilizes Principal Component Analysis (PCA) to identify the
key features that capture the most variance in the training images.

Given a set of training images $I_1, I_2, \ldots, I_n$, the average face is computed as:
$\bar{I} = \frac{1}{n} \sum_{i=1}^n I_i$. For each image, compute the mean-subtracted image
$\Phi_i = I_i - \bar{I}$ so that the data is centered around the average face needed for PCA.
The covariance matrix $C$ is then computed as

\begin{equation}
	C = AA^T = \frac{1}{n} \sum_{i=1}^n \Phi_i \Phi_i^T
\end{equation}

where $A = [\Phi_1, \Phi_2, \ldots, \Phi_n]$. For PCA, the eigenvectors and eigenvalues must be
computed but the operation $AA^T$ is large and intractable. Instead, it can be shown that the
$N^2$ eigenvectors $v_l$ from $A^TA$ correspond to the largest from $AA^T$ with
respect to the eigenvalues. Once $v_l$ is computed, it must be transformed back to the original
space to get the eigenvectors $u_l$
\begin{equation}
	u_l = \sum{k=1}^N v_{lk} \Phi_k
\end{equation}

The eigenvectors of $C$, $u_l$, are the eigenfaces, which form the basis of the face space.
Not all eigenfaces are needed to represent face space, so the top $K$ principal components
that retain a certain percentage of information from the original data.

\subsection{Image Reconstruction}

This section covers how images are represented in the new face space and how they can be
reconstructed from the eigen-coefficients.

Eigen-coefficients can be calculated by projecting the mean-subtracted images $\Phi_i$ onto
the $K$ eigenfaces to compute the eign-coefficients $\omega_k$ as follows:
\begin{equation}
	\omega_k = u_k^T \Phi_i
\end{equation}
And $\Omega^T = [\omega_1, \omega_2, \ldots, \omega_K]$ can be stacked to represent the image
in the new face space.

The original image can be reconstructed from the eigen-coefficients and the average face as:
\begin{equation}
	\hat{I} = \sum_{k=1}^K \omega_k u_k + \bar{I}
\end{equation}

A given image can now be compared to the reconstructed image to determine how well it is
represented in the face space. The distance between the original image and the reconstructed
image can be computed using the Euclidean distance:
\begin{equation}
	d_E = \sqrt{\sum_{i=1}^N (I_i - \hat{I}_i)^2}
\end{equation}
\label{eq:euclidean}
where $N$ is the number of pixels in the image. A smaller distance indicates a better
reconstruction, while a larger distance indicates that the image is not well represented in the
face space.

\subsection{Experiment (a)}

This section outlines how query images can be matched against gallery images to evaluate
the performance of the face recognition system.

The closest match of a given query image $q_i$ is found by computing the distance to all
galley images $g_j$ in the face space and selecting the one with the smallest distance.
The distance can be computed using the Euclidean distance \ref{eq:euclidean}
or the Mahalanobis distance:

\begin{equation}
	d_M = \sqrt{(\Omega - \Omega_i)^T C^{-1} (\Omega - \Omega_i)}
\end{equation}
where $C$ is the covariance matrix of the eigen-coefficients. The Mahalanobis distance
accounts for the variance in the data and can provide better discrimination between different
faces. When images are projected to face space, the scale of the eigen-coefficients can vary
significantly. If the Euclidean distance is used, it may be dominated by larger components that
capture irrelevant variations (i.e. lighting). The Mahalanobis fixes this by accounting for the
distribution of the data and scales the coordinates according the covariance matrix.

To evaluate the performance of the facial recognition system, the Cumulative Match Characteristic
(CMC) curve is used. The CMC curve plots the accuracy, or recognition rate, as the rank varies.
The rank is the acceptable range in which a correct match can be found (i.e. rank 1 means the
correct match must be the closest).

\subsection{Experiment (b)}

To identify intruders, a threshold can be defined on the distance between the query image and
the closest gallery image in the face space. If the distance exceeds the threshold, the query
image is classified as an intruder. The choice of threshold is critical, as it affects the
trade-off between false positives (incorrectly classifying a legitimate user as an intruder)
and false negatives (failing to identify an intruder). To evaluate the performance of the
intruder detection system, the Receiver Operating Characteristic (ROC) is used. The ROC curve
plots the true positive rate against the false positive rate as the threshold varies. The area
under the ROC curve (AUC) can be used as a summary measure of the system's performance, with
a higher AUC indicating better performance.

\subsection{Experiment (c)}

Varying the resolution of the input images can have significant effects on the performance
of the face recognition system. Higher resolution images contain more detailed information,
which can provide more discriminant features for recognition.

% ============================================================
%  SECTION 2 -- RESULTS AND DISCUSSION
% ============================================================
\section{Results and Discussion}

% Present and discuss your results. Include sample results, discuss major
% findings, and reference any figures or tables generated (i.e., using figure/table captions) to illustrate
% your results. For clarity, the results for each part must be presented in a separate subsection.

% ── Experiment 1 ────────────────────────────────────────────
\subsection{Experiment (a)}

\subsubsection{a.I}

\subsubsection{a.II}

\subsubsection{a.III}

\subsubsection{a.IV}

\subsubsection{a.V}

\subsection{Experiment (b)}

\subsection{Experiment (c)}

\subsection{Experiment (d)}

\subsection{Experiment (e)}

\end{document}
